\documentclass{article}
\usepackage[utf8]{inputenc}
\usepackage{hyperref}

\title{Notes on Reasoning in Large Language Models}
\author{Ana Souza}
\date{}

\begin{document}
\maketitle

\begin{abstract}
We collect what is known about eliciting multi-step reasoning from large language models
and how that behaviour is steered after pre-training.
\end{abstract}

\section{Background}
Open foundation models such as LLaMA~\cite{touvron2023llama} make it possible to study
prompting methods without access to a proprietary system.

\section{Prompting for reasoning}
Chain-of-thought prompting~\cite{wei2022chainofthought} shows that intermediate steps
improve accuracy on arithmetic and commonsense benchmarks. Sampling several reasoning paths
and voting on the answer~\cite{wang2022selfconsitency} improves on a single greedy path, and
interleaving reasoning with tool calls~\cite{yao2022react} lets a model act on what it reads.
Recent work reports that these gains transfer to biomedical text~\cite{chen2023benchmarking}.

\section{Steering behaviour}
Preference data can replace an explicit reward model~\cite{rafailov2023direct}, and written
principles can replace most human harm labels~\cite{bai2022constitutional}.

\bibliographystyle{plain}
\bibliography{labshelf,refs}

\end{document}
